% common/pool/立体几何/难题/q_solid_frustum_angle_arc_length.tex
% ★★★ 难题

\newcommand{\qSolidFrustumAngleArcLengthStem}{
已知正三棱台 $ABC-A_1B_1C_1$ 的上、下底面边长分别为 $4$ 和 $6$，侧棱长为 $2$。点 $P$ 在侧面 $BCC_1B_1$ 内运动（包含边界），且 $AP$ 与平面 $BCC_1B_1$ 所成角的正切值为 $\sqrt{6}$。求所有满足条件的动点 $P$ 形成的轨迹长度。
}

\newcommand{\qSolidFrustumAngleArcLengthOptions}{
A. $\frac{4\pi}{3}$ \quad
B. $\frac{\sqrt{5}\pi}{3}$ \quad
C. $\frac{\sqrt{7}\pi}{3}$ \quad
D. $\frac{2\pi}{3}$
}

\newcommand{\qSolidFrustumAngleArcLengthFigure}[1][1]{
\begin{tikzpicture}[scale=#1, line join=round, line cap=round, >=stealth]
    % 定义底部三角形顶点
    \coordinate (B) at (0, 0);
    \coordinate (C) at (6, 0);
    \coordinate (A) at (3, 1.8);
    
    % 定义上部三角形顶点
    \coordinate (B1) at (1, 2.5);
    \coordinate (C1) at (5, 2.5);
    \coordinate (A1) at (3, 3.7);

    % 连线
    \draw[dashed, thin, gray] (A) -- (B);
    \draw[dashed, thin, gray] (A) -- (C);
    \draw[thick] (B) -- (C);
    \draw[thick] (B) -- (B1);
    \draw[thick] (C) -- (C1);
    \draw[thick] (B1) -- (C1);
    \draw[dashed, thin, gray] (A) -- (A1);
    \draw[thick] (A1) -- (B1);
    \draw[thick] (A1) -- (C1);

    % 标记顶点
    \node[below left] at (B) {$B$};
    \node[below right] at (C) {$C$};
    \node[above right] at (C1) {$C_1$};
    \node[above left] at (B1) {$B_1$};
    \node[above] at (A1) {$A_1$};
    \node[above right] at (A) {$A$};

    % 侧面 BCC1B1 上的圆弧轨迹
    \coordinate (H) at (3, 2.5);
    \coordinate (P1) at (2, 0);
    \coordinate (P2) at (4, 0);
    
    \draw[thick, densely dashed] (B1) to[bend right=20] (P1);
    \draw[thick, densely dashed] (C1) to[bend left=20] (P2);
    
    \node[above] at (H) {$H$};
    \fill (H) circle (1.2pt);
    \node[below] at (P1) {$P_1$};
    \node[below] at (P2) {$P_2$};
\end{tikzpicture}
}

\newcommand{\qSolidFrustumAngleArcLengthAnswer}{
A
}

\newcommand{\qSolidFrustumAngleArcLengthSolution}{
\textbf{第一步：还原为正四面体} \\
将正三棱台 $ABC-A_1B_1C_1$ 延长各侧棱，使其交于顶点 $V$。
由于上下底面边长比为 $B_1C_1 : BC = 4 : 6 = 2 : 3$，利用相似比可知 $VB_1 : VB = 2 : 3$。
又因截去的侧棱长 $BB_1 = 2$，所以 $VB - VB_1 = 2$，解得 $VB = 6, VB_1 = 4$。
由此可见，补全后的三棱锥 $V-ABC$ 其侧棱与底面边长均为 $6$，即它是一个棱长为 $6$ 的正四面体。

\textbf{第二步：确定点 $A$ 在平面 $BCC_1B_1$ 的投影位置} \\
平面 $BCC_1B_1$ 即为正四面体的侧面 $VBC$。
在棱长为 $6$ 的正四面体中，顶点 $A$ 到对立面 $VBC$ 的距离（高）为：
$h = \frac{\sqrt{6}}{3} \times 6 = 2\sqrt{6}$。
设 $A$ 在平面 $VBC$ 上的投影为 $H$，则 $H$ 是正三角形 $VBC$ 的中心。
$\triangle VBC$ 的高为 $\frac{\sqrt{3}}{2} \times 6 = 3\sqrt{3}$，中心 $H$ 到顶点 $V$ 的距离为 $2\sqrt{3}$，到 $BC$ 的距离为 $\sqrt{3}$。
而线段 $B_1C_1$ 在侧面上与 $BC$ 平行且距离也恰好为 $\sqrt{3}$（因为 $VB_1 = 4$ 对应的高度比例），因此点 $H$ 恰好位于线段 $B_1C_1$ 的中点上。

\textbf{第三步：分析线面角条件，转化为圆轨迹} \\
根据题意，$AP$ 与平面 $BCC_1B_1$ 所成角的正切值为 $\sqrt{6}$。
在直角 $\triangle AHP$ 中，$AH \perp \text{平面 } BCC_1B_1$，故 $\angle APH$ 即为 $AP$ 与该平面所成的角。
由 $\tan \angle APH = \frac{AH}{HP} = \frac{2\sqrt{6}}{HP} = \sqrt{6}$，解得 $HP = 2$。
这说明点 $P$ 的轨迹是平面 $BCC_1B_1$ 内，以 $H$ 为圆心，半径为 $2$ 的圆。

\textbf{第四步：求梯形区域内轨迹的有效长度} \\
点 $P$ 必须在梯形 $BCC_1B_1$ 内部或边界上运动。
因 $H$ 是 $B_1C_1$（长为 $4$）的中点，所以 $HB_1 = HC_1 = 2$。即圆弧恰好经过点 $B_1$ 和 $C_1$。
以 $H$ 为原点，在梯形平面内建立直角坐标系。梯形区域纵坐标范围为 $y \in [-\sqrt{3}, 0]$。
圆的方程为 $x^2 + y^2 = 4$。
当 $y = -\sqrt{3}$（即圆与底边 $BC$ 相交时），代入得 $x^2 + 3 = 4 \implies x = \pm 1$。
圆心对应的交点角度满足 $\sin \theta = -\frac{\sqrt{3}}{2}$，即角度为 $240^\circ$ 和 $300^\circ$。
满足梯形范围（$-\sqrt{3} \leqslant y \leqslant 0$）的有效圆弧对应角度区间为 $[180^\circ, 240^\circ]$ 与 $[300^\circ, 360^\circ]$。
这部分轨迹的圆心角总计为 $(240^\circ - 180^\circ) + (360^\circ - 300^\circ) = 120^\circ = \frac{2\pi}{3}$ 弧度。
所以，动点 $P$ 的有效轨迹长度为 $L = \alpha \times r = \frac{2\pi}{3} \times 2 = \frac{4\pi}{3}$。
故选 A。
}

\newcommand{\qSolidFrustumAngleArcLengthDiagnosis}{
\textbf{学生可能存在的问题：}
1. \textbf{模型还原意识弱}：面对棱台，没有第一时间想到补全为完整的棱锥，无法利用正四面体“定型定性”的优美性质来简化计算。
2. \textbf{空间投影位置找不准}：虽然知道求高，但无法通过几何比例发现投影点 $H$ 恰好在 $B_1C_1$ 的中点，导致无法精准锁定圆心的位置。
3. \textbf{线面角定义不清}：不知道线面角需在射影直角三角形 $\triangle AHP$ 中求解，导致无法将正切值条件转化为动点到定点 $H$ 的距离为定值。
4. \textbf{忽视边界截断}：算出了 $HP=2$ 的半径后，直接算整个圆的周长或半圆周长，没有代数建系去严谨计算圆弧与梯形 $BC$ 边界的截断角度。
}

\newcommand{\qSolidFrustumAngleArcLengthTeachingNotes}{
\textbf{课堂教学提示：}
1. \textbf{还原模型（遇台先补锥）}：由于是正三棱台，补全后必然是正三棱锥。通过边长比 $4:6$ 和侧棱 $2$ 快速算出侧棱为 $6$，发现这是个正四面体，直接破题。
2. \textbf{空间降维打击}：空间中关于线面角的条件 $\tan \theta = \sqrt{6}$，本质上是给出“点面距离”与“面内射影距离”之比。引导学生将其剥离出来，降维成纯平面的圆轨迹截断问题。
3. \textbf{精确画图与代数验证}：在确定梯形面内的圆弧有效段时，切忌凭感觉画图。引导学生在梯形面内建系，利用直线 $y = -\sqrt{3}$ 截圆 $x^2 + y^2 = 4$，用代数精确算出交点坐标，从而严谨得出 $120^\circ$ 的圆心角。
}
